\documentclass[11pt]{article}

\usepackage[a4paper,margin=1in]{geometry}
\usepackage{amsmath,amssymb,bm,mathtools}
\usepackage{microtype}
\usepackage{xcolor}
\usepackage[most]{tcolorbox}

\definecolor{HessianPurple}{HTML}{6F42C1}
\definecolor{HessianPurpleLight}{HTML}{F5F1FC}
\newtcolorbox{hessianhighlight}{
  enhanced,
  breakable,
  colback=HessianPurpleLight,
  colframe=HessianPurple,
  colbacktitle=HessianPurple,
  coltitle=white,
  fonttitle=\bfseries,
  title={Distributional dependence of the Hessian blocks},
  boxrule=0.8pt,
  arc=2pt,
  left=7pt,
  right=7pt,
  top=6pt,
  bottom=6pt
}
\usepackage[
  colorlinks=true,
  linkcolor=HessianPurple,
  citecolor=HessianPurple,
  urlcolor=HessianPurple
]{hyperref}

\newcommand{\R}{\mathbb{R}}
\newcommand{\E}{\mathbb{E}}
\newcommand{\tr}{\operatorname{tr}}
\newcommand{\diag}{\operatorname{diag}}
\newcommand{\BlockDiag}{\operatorname{BlockDiag}}
\newcommand{\ip}[2]{\left\langle #1,#2\right\rangle_F}

\title{Hessian Structure of a Linear Bigram Model}
\author{Senmiao Wang}
\date{}

\begin{document}
\maketitle

\begin{abstract}
This note derives the gradient and Hessian of an untied two-layer linear
bigram model trained with the full-vector mean squared error.  The diagonal
Hessian blocks reveal two distinct effects of data imbalance.  Across
embedding columns, block heterogeneity is determined solely by the input-token
distribution.  Across rows of the language-model head, all Hessian blocks are
identical for arbitrary input and output distributions: the input distribution
changes their common value, whereas the conditional target distribution does
not enter these blocks.  The conditional target distribution affects the
gradient and the residual-dependent part of the cross-layer Hessian.
\end{abstract}

\section{Problem Setting}

We study population risk minimization for bigram prediction.  We first specify
the data-generating distribution, which determines the form of token
imbalance, and then introduce the two-layer linear network and its training
objective.

\paragraph{Bigram prediction problem and data distribution.}
Let \(C\) be the number of token classes and let
\(\bm{e}_1,\ldots,\bm{e}_C\in\R^C\) be the standard basis vectors.  A bigram
sample consists of an input token \(x=\bm{e}_i\) and its next-token target
\(y=\bm{e}_j\).  We assume that \((x,y)\) follows a population distribution
\(\mathcal{D}\) specified by the input-token probabilities
\(\pi=(\pi_1,\ldots,\pi_C)\) and the conditional next-token probabilities
\begin{equation*}
\pi_i=\Pr(x=\bm{e}_i),
\qquad
\pi_{j\mid i}=\Pr(y=\bm{e}_j\mid x=\bm{e}_i),
\end{equation*}
where \(\sum_i\pi_i=1\) and \(\sum_j\pi_{j\mid i}=1\) for every input class
with \(\pi_i>0\).  The joint and output-token probabilities are therefore
\begin{equation*}
\Pr(x=\bm{e}_i,y=\bm{e}_j)=\pi_i\pi_{j\mid i},
\qquad
\pi_j^{\mathrm{out}}
=\Pr(y=\bm{e}_j)
=\sum_{i=1}^C\pi_i\pi_{j\mid i}.
\end{equation*}
For each input class, collect its conditional next-token distribution in
\begin{equation*}
\bm{\pi}_{\cdot\mid i}
=
\begin{bmatrix}
\pi_{1\mid i}&\cdots&\pi_{C\mid i}
\end{bmatrix}^{\!\top}
=\E[y\mid x=\bm{e}_i].
\end{equation*}
We also define the diagonal matrix of input-token probabilities as
\begin{equation*}
D_\pi=\diag(\pi_1,\ldots,\pi_C).
\end{equation*}
In this formulation, input-token imbalance corresponds to nonuniform
\(\pi_i\), while the conditional probabilities \(\pi_{j\mid i}\) determine
how the next-token distribution varies with the input.

\paragraph{Two-layer linear network and training objective.}
Let \(d\) be the hidden dimension.  The network consists of an embedding
matrix \(W\) and a language-model head \(V\):
\begin{equation*}
W\in\R^{d\times C},
\qquad
V\in\R^{C\times d}.
\end{equation*}
For an input \(x=\bm{e}_i\), the \(i\)-th embedding column
\(w_i=W\bm{e}_i\) is the hidden representation, and the network prediction is
\begin{equation*}
h=Wx=w_i,
\qquad
f(V,W;x)=Vh=VWx.
\end{equation*}
Given a target \(y\), define the prediction error by
\begin{equation*}
e=f(V,W;x)-y=VWx-y.
\end{equation*}
Here \(e\) denotes the prediction error, whereas the bold symbol
\(\bm{e}_i\) denotes a standard basis vector.

We train both layers by minimizing the population full-vector MSE:
\begin{equation}
\mathcal{L}(V,W)
=\frac12\E_{(x,y)\sim\mathcal{D}}\!\left[\|VWx-y\|_2^2\right]
=\frac12\E_{(x,y)\sim\mathcal{D}}\!\left[e^\top e\right].
\label{eq:loss}
\end{equation}
Conditioning \eqref{eq:loss} on the input token yields
\begin{equation*}
\mathcal{L}(V,W)
=\frac12\sum_{i=1}^C
\pi_i\|Vw_i-\bm{\pi}_{\cdot\mid i}\|_2^2
+\text{constant}.
\end{equation*}
Thus, up to a parameter-independent term, training is a weighted
factorization of the conditional bigram probabilities, with weights given by
the input-token distribution.

\section{Calculation of the Gradients and Hessian}

Because \(x\) and \(y\) are fixed under parameter differentiation, the first
differential of the error is
\begin{equation*}
de=d(VWx)=dV\,Wx+V\,dW\,x.
\end{equation*}
It follows from \eqref{eq:loss} that
\begin{equation*}
d\mathcal{L}
=\E\!\left[
e^\top\!\left(dV\,Wx+V\,dW\,x\right)
\right].
\end{equation*}
Using the Frobenius inner product
\(\ip{A}{B}=\tr(A^\top B)\), the two terms satisfy
\begin{equation*}
e^\top dV\,h=\ip{eh^\top}{dV},
\qquad
e^\top V\,dW\,x=\ip{V^\top e x^\top}{dW}.
\end{equation*}
Therefore, the population gradients are
\begin{equation}
\begin{aligned}
\nabla_V\mathcal{L}
&=\E[eh^\top]
=\sum_{i=1}^C
\pi_i\left(Vw_i-\bm{\pi}_{\cdot\mid i}\right)w_i^\top,\\
\nabla_W\mathcal{L}
&=\E[V^\top e x^\top]
=\sum_{i=1}^C
\pi_iV^\top\left(Vw_i-\bm{\pi}_{\cdot\mid i}\right)\bm{e}_i^\top.
\end{aligned}
\label{eq:gradients}
\end{equation}
In particular, the gradient with respect to the \(i\)-th embedding column is
\begin{equation*}
\nabla_{w_i}\mathcal{L}
=\pi_iV^\top\left(Vw_i-\bm{\pi}_{\cdot\mid i}\right).
\end{equation*}
Equation~\eqref{eq:gradients} shows that both the input probability \(\pi_i\)
and the conditional target probabilities \(\pi_{j\mid i}\) affect the
gradients.

The model output is bilinear in \(V\) and \(W\), so differentiating the error
a second time gives
\begin{equation*}
d^2e=d^2(VWx)=2\,dV\,dW\,x.
\end{equation*}
The factor of two collects the two product-rule contributions.  Hence the
complete Hessian quadratic form is
\begin{equation}
d^2\mathcal{L}
=\E\!\left[
\left\|dV\,Wx+V\,dW\,x\right\|_2^2
+2e^\top dV\,dW\,x
\right].
\label{eq:full-hessian}
\end{equation}
Equivalently,
\begin{equation*}
\begin{aligned}
d^2\mathcal{L}
=\E\Big[
&\|dV\,Wx\|_2^2
+\|V\,dW\,x\|_2^2\\
&+2(dV\,Wx)^\top V\,dW\,x
+2e^\top dV\,dW\,x
\Big].
\end{aligned}
\end{equation*}
The first two terms define the two within-layer, or diagonal, Hessian blocks.
The final two terms define the cross-layer blocks.  The prediction error
appears only in the cross-layer contribution because the model is linear in
each parameter matrix separately but bilinear in the pair.

\section{Hessian Block Structure and Effect of Data Imbalance}
\label{sec:hessian-blocks}

We now state the Hessian block structure and its distributional dependence.
The detailed blockwise calculations are deferred to
Appendix~\ref{app:hessian-calculations}.

\paragraph{Embedding-layer diagonal block.}
The embedding matrix is partitioned by columns along the input-vocabulary
direction,
\begin{equation*}
W=
\begin{bmatrix}
w_1&\cdots&w_C
\end{bmatrix},
\end{equation*}
with one parameter block \(w_i\in\R^d\) for each input token.  Under the
corresponding column-wise vectorization, the embedding-layer Hessian is
\begin{equation}
\begin{aligned}
H_{WW}
&=\BlockDiag\!\left(
\pi_1V^\top V,\ldots,\pi_CV^\top V
\right)
=D_\pi\otimes V^\top V,\\
H_{WW}^{(i)}
&=\pi_iV^\top V,
\qquad
H_{WW}[dW]=V^\top V\,dW\,D_\pi.
\end{aligned}
\label{eq:embedding-hessian}
\end{equation}
At fixed \(V\), all embedding blocks have the common shape \(V^\top V\);
their block-to-block heterogeneity is determined entirely by the input
probabilities \(\pi_i\).  Neither the conditional target distribution nor the
output marginal appears in \eqref{eq:embedding-hessian}.
Equivalently, its input-vocabulary block structure is
\begin{equation*}
H_{WW}
=
\begin{bmatrix}
\pi_1V^\top V & 0 & \cdots & 0\\
0 & \pi_2V^\top V & \cdots & 0\\
\vdots & \vdots & \ddots & \vdots\\
0 & 0 & \cdots & \pi_CV^\top V
\end{bmatrix}.
\end{equation*}

\paragraph{LM-head diagonal block.}
The LM head is partitioned by rows along the output-vocabulary direction,
\begin{equation*}
V=
\begin{bmatrix}
v_1^\top\\
\vdots\\
v_C^\top
\end{bmatrix},
\end{equation*}
with one parameter block \(v_a\in\R^d\) for each output token.  The LM-head
Hessian satisfies
\begin{equation}
\begin{aligned}
H_{VV}
&=\BlockDiag\!\left(
WD_\pi W^\top,\ldots,WD_\pi W^\top
\right)
=I_C\otimes\left(WD_\pi W^\top\right),\\
H_{VV}^{(1)}
&=\cdots=H_{VV}^{(C)}=WD_\pi W^\top,\\
H_{VV}[dV]
&=dV\,WD_\pi W^\top.
\end{aligned}
\label{eq:head-hessian}
\end{equation}
Thus all output-class blocks are identical for arbitrary input and output
distributions.  Their common block is
\(WD_\pi W^\top=\sum_i\pi_iw_iw_i^\top\).  Therefore, nonuniform input
probabilities can unequally weight the embedding directions and induce or
amplify spectral anisotropy within every common block, even though they cannot
create block-to-block heterogeneity across output tokens.  The target
distribution does not enter \eqref{eq:head-hessian}.
Equivalently, its output-vocabulary block structure is
\begin{equation*}
H_{VV}
=
\begin{bmatrix}
WD_\pi W^\top & 0 & \cdots & 0\\
0 & WD_\pi W^\top & \cdots & 0\\
\vdots & \vdots & \ddots & \vdots\\
0 & 0 & \cdots & WD_\pi W^\top
\end{bmatrix}.
\end{equation*}

\paragraph{Cross-layer blocks.}
Let \(v_a^\top\) be the \(a\)-th row of \(V\), and define the conditional mean
error
\begin{equation*}
\bar e_{i,a}=v_a^\top w_i-\pi_{a\mid i}.
\end{equation*}
With the convention that the Hessian quadratic form contains
\(2\,dv_a^\top H_{v_a,w_i}dw_i\), the cross-layer block is
\begin{equation}
H_{v_a,w_i}
=\pi_i\left(w_iv_a^\top+\bar e_{i,a}I_d\right).
\label{eq:cross-block}
\end{equation}
Unlike the two diagonal Hessian blocks, the cross-layer block depends on the
conditional target distribution through \(\bar e_{i,a}\).  When
\(Vw_i=\bm{\pi}_{\cdot\mid i}\), this residual-dependent term vanishes, while
the Gauss--Newton coupling \(\pi_iw_iv_a^\top\) generally remains.

\begin{hessianhighlight}
At fixed network parameters \(W\) and \(V\):
\begin{itemize}
  \item The heterogeneity of the embedding-layer blocks depends directly
  only on the input-token probabilities \(\pi_i\).  It is independent of
  \(\pi_{j\mid i}\) and the output marginal.
  \item All LM-head blocks are identical and equal to \(WD_\pi W^\top\).
  Input imbalance cannot create block-to-block heterogeneity across output
  tokens, but its unequal weighting of embedding directions can induce or
  amplify spectral heterogeneity within every common block.
  \item The conditional target distribution affects Hessian curvature only
  through the residual-dependent part of the cross-layer blocks.
\end{itemize}
\end{hessianhighlight}

\newpage
\appendix
\section{Detailed Hessian Block Calculations}
\label{app:hessian-calculations}

This appendix derives the block formulas stated in
Section~\ref{sec:hessian-blocks}.  We start from the complete Hessian
quadratic form in \eqref{eq:full-hessian}.

\paragraph{Embedding-layer diagonal block.}
Setting \(dV=0\) in \eqref{eq:full-hessian} gives
\begin{equation*}
d_W^2\mathcal{L}
=\E\!\left[\|V\,dW\,x\|_2^2\right].
\end{equation*}
Since \(x=\bm{e}_i\), write \(dw_i=dW\bm{e}_i\).  Averaging over the input
distribution gives
\begin{equation*}
d_W^2\mathcal{L}
=\sum_{i=1}^C
\pi_i\,dw_i^\top V^\top V\,dw_i.
\end{equation*}
There are no terms coupling \(dw_i\) and \(dw_k\) for \(i\neq k\), which
yields \eqref{eq:embedding-hessian}.  Consequently, for any homogeneous
matrix norm and any \(\pi_k>0\),
\begin{equation*}
\frac{\|H_{WW}^{(i)}\|}{\|H_{WW}^{(k)}\|}
=\frac{\pi_i}{\pi_k}.
\end{equation*}
If \(\pi_i=0\), then \(H_{WW}^{(i)}=0\), so the corresponding embedding is
unidentified by the observed bigram distribution.  If every \(\pi_i>0\) and
\(V\) has full column rank, then
\begin{equation*}
\kappa(H_{WW})
=\frac{\pi_{\max}}{\pi_{\min}}\kappa(V^\top V)
=\frac{\pi_{\max}}{\pi_{\min}}\kappa(V)^2.
\end{equation*}
Thus input imbalance and the geometry of the LM head contribute
multiplicatively to ill-conditioning in the embedding layer.

\paragraph{LM-head diagonal block.}
Setting \(dW=0\) in \eqref{eq:full-hessian} gives
\begin{equation*}
d_V^2\mathcal{L}
=\E\!\left[\|dV\,Wx\|_2^2\right].
\end{equation*}
Writing the \(a\)-th row of \(dV\) as \(dv_a^\top\), we obtain
\begin{equation*}
d_V^2\mathcal{L}
=\tr\!\left(dV^\top dV\,WD_\pi W^\top\right)
=\sum_{a=1}^C
dv_a^\top\left(WD_\pi W^\top\right)dv_a,
\end{equation*}
which yields \eqref{eq:head-hessian}.  The target independence follows
directly from
\begin{equation*}
d_V^2\!\left[\frac12\|Vw_i-y\|_2^2\right]
=\|dV\,w_i\|_2^2,
\end{equation*}
which does not depend on \(y\).

\paragraph{Cross-layer blocks.}
The cross-layer contribution in \eqref{eq:full-hessian} is
\begin{equation*}
\begin{aligned}
d_{VW}^2\mathcal{L}
=2\sum_{i=1}^C\pi_i\Big[
&(dV\,w_i)^\top V\,dw_i\\
&+\left(Vw_i-\bm{\pi}_{\cdot\mid i}\right)^\top dV\,dw_i
\Big].
\end{aligned}
\end{equation*}
The bilinear term between \(dv_a\) and \(dw_i\) is
\begin{equation*}
2\pi_i\,dv_a^\top
\left(w_iv_a^\top+\bar e_{i,a}I_d\right)dw_i.
\end{equation*}
Under the block convention stated before \eqref{eq:cross-block}, this gives
the cross-layer result in \eqref{eq:cross-block}.

\end{document}
